\section{Sparse Distributed Memory: cómo escribir y cómo leer a partir de consultas con ruido}

La primera parte de la clase todavía no trata el Transformer; construye un sistema de memoria asociativa. El objetivo no es leer un registro con una dirección exacta, sino que una query cercana siga recuperando el pattern correcto, y obtener al mismo tiempo alta capacidad, tolerancia a fallos y almacenamiento distribuido. La técnica central de SDM es fijar un conjunto de hard locations dispersas, de modo que cada escritura y cada lectura actúen sobre un grupo de posiciones dentro del radio de la query.

\subsection{Objetivos de diseño: capacidad, robustez y fault tolerance}

La diapositiva de requisitos enumera cuatro: high memory capacity, robustness to query noise, biological plausibility y fault tolerance. La siguiente diapositiva separa sparse de distributed: el espacio de direcciones tiene una dimensión altísima y las neuronas reales ocupan solo una fracción mínima de las posiciones posibles; sin embargo, un pattern activa a la vez varias posiciones cercanas, por lo que ninguna dirección individual contiene la memoria completa.

\lecturefigure{slide-04-brain-memory-requirements.jpg}{Requisitos de diseño de una memoria asociativa de tipo cerebral}{Diapositiva recuperada del video oficial V004; 00:02:12}

\lecturefigure{slide-05-sdm-unique.jpg}{Los dos sentidos de SDM: sparse y distributed}{Diapositiva recuperada del video oficial V005; 00:02:48}

\begin{knowledgebox}{Glosario: address, pattern, value y hard location}
\term{pattern} es el vector de alta dimensión que se escribe o que se usa como query; \term{hard location} es una dirección dispersa predefinida por el sistema; \term{address/key} determina qué posiciones se activan; \term{value} es el contenido que se escribe y que se devuelve en la consulta. En la auto-associative memory clásica, key y value pueden ser iguales; en la hetero-associative memory pueden ser distintos.
\end{knowledgebox}

Sea el espacio binario de direcciones $\{0,1\}^{n}$, la $i$-ésima hard location $h_i$, la dirección de escritura $p$ y el radio $d$. El indicador de activación es
$$
a_i(p)=\mathbf{1}\!\left[D_H(h_i,p)\le d\right],
$$
donde $D_H$ es la Hamming distance, $n$ es la dimensión de los vectores y $d$ controla cuántas posiciones se activan. En alta dimensión, la mayoría de los vectores aleatorios están muy lejos entre sí, y aun así un radio local puede contener suficientes hard locations para realizar un cómputo distribuido.

\subsection{Escritura: posiciones cercanas y superposition}

En la primera figura de escritura, el pattern verde activa las hard locations dentro de su radio; en la segunda se agrega el pattern naranja, y los radios de ambas escrituras pueden traslaparse; en la tercera se agrega el pattern azul, y las posiciones del traslape guardan a la vez las contribuciones de varios value. La superposition no concatena varias memorias: las acumula sobre el mismo vector de almacenamiento, y en la lectura el ruido se elimina mediante la agregación estadística de muchas posiciones.

\lecturefigure{slide-06-sdm-write-pattern1.jpg}{SDM escribe el primer pattern: se activan las hard locations dentro del radio}{Diapositiva recuperada del video oficial V006; 00:03:26}

\lecturefigure{slide-07-sdm-write-pattern2.jpg}{Escritura del segundo pattern: los vecindarios de las direcciones pueden traslaparse}{Diapositiva recuperada del video oficial V007; 00:03:42}

\lecturefigure{slide-08-sdm-write-superposition.jpg}{Varios pattern se almacenan superpuestos en las hard locations}{Diapositiva recuperada del video oficial V008; 00:04:14}

Si la $i$-ésima posición guarda un vector de contadores $w_i$, escribir el value $v$ en la dirección $p$ puede expresarse como
$$
w_i\leftarrow w_i+a_i(p)\,v.
$$
Aquí $v$ puede ser un vector bipolar o un vector de valores reales, y $a_i(p)$ determina si la posición participa. Una posición acumula varios $v$, y la capacidad y el ruido provienen del mismo mecanismo: cuantos más pattern comparten posiciones, mayor es el aprovechamiento del almacenamiento y también mayor la interferencia cruzada.

\begin{lstlisting}[caption={Escritura distribuida de SDM (pseudocódigo conceptual)}]
def sdm_write(hard_locations, counters, address, value, radius):
    for index, location in enumerate(hard_locations):
        if hamming(location, address) <= radius:
            counters[index] += value
\end{lstlisting}

\teachervoice{El ponente recuerda varias veces al público que primero piense en estas «neuron» como hard locations abstractas y que solo después se relacionarán con la biología. Este orden evita interpretar un diagrama bidimensional de círculos como un espacio real dentro del cerebro: la distancia en la figura representa una Hamming distance de alta dimensión, y los círculos son solo una proyección didáctica.}

\subsection{Lectura: de las neuronas activadas a la intersección de pattern}

La operación de lectura tiene dos perspectivas equivalentes. En la neuron view, la query activa las posiciones cercanas, se suman los vectores almacenados en cada posición y luego se aplica majority vote dimensión por dimensión; la pattern view abstrae las hard locations y examina directamente el tamaño de la intersección entre la query sphere y cada write sphere. Cuanto mayor es la intersección, más posiciones se activaron tanto al escribir ese pattern como al leer con la query actual, y mayor es el peso del value de ese pattern.

\lecturefigure{slide-09-sdm-read-neurons.jpg}{Neuron view: una query con ruido activa posiciones cercanas y agrega lo almacenado}{Diapositiva recuperada del video oficial V009; 00:05:02}

\lecturefigure{slide-10-sdm-read-patterns.jpg}{Pattern view: la sphere intersection determina el peso de cada memoria}{Diapositiva recuperada del video oficial V010; 00:05:40}

\lecturefigure{slide-11-circle-intersection.jpg}{Intersección y distancia entre dos esferas de alta dimensión}{Diapositiva recuperada del video oficial V011; 00:05:56}

La forma continua de la lectura puede escribirse como
$$
r(q)=\sum_i a_i(q)w_i,
\qquad
\hat v_j=g(r_j)=\mathbf{1}[r_j>0],
$$
donde $q$ es la query, $r(q)$ es la lectura total de las posiciones activadas y $g$ es un umbral por dimensión o un majority decoder. Al sustituir la fórmula de escritura, la contribución de cada pattern $p_\mu$ a la salida es proporcional a
$$
I\!\left(D_H(p_\mu,q);d,n\right)
=\left|B(p_\mu,d)\cap B(q,d)\right|,
$$
es decir, al tamaño de la intersección de dos Hamming ball; $B(p,d)$ denota la esfera con centro en $p$ y radio $d$.

\lecturefigure{slide-12-query-read-equation.jpg}{Las cuatro etapas de la fórmula de query read: intersección, ponderación, normalización y decodificación}{Diapositiva recuperada del video oficial V012; 00:07:18}

\begin{knowledgebox}{Lectura de la figura: por qué una noisy query aún puede recuperar el pattern}
La query no tiene que ser igual a la dirección original: mientras siga traslapándose ampliamente con la sphere del pattern correcto, activará muchas de las posiciones donde se escribió ese value. Cuando errores aleatorios alteran unos cuantos bit, la intersección correcta suele seguir siendo mayor que la de los pattern no relacionados; el majority vote promedia el ruido independiente. La recuperación falla cuando los pattern son demasiado densos, el radio es demasiado grande o la query está demasiado lejos del pattern original.
\end{knowledgebox}

\subsection{Resumen de la sección}

SDM usa direcciones dispersas de alta dimensión para realizar escrituras distribuidas y la sphere intersection para el direccionamiento por contenido. El radio de escritura controla a la vez la capacidad y la interferencia; la lectura elimina el ruido mediante la agregación de las posiciones activadas y la majority decoding. En la siguiente sección basta demostrar que el intersection weight tiene una forma exponencial similar a softmax para conectar estas operaciones de memoria con attention.
